\documentclass[11pt,a4paper,twoside]{article}
\usepackage[utf8]{inputenc}
\usepackage[english]{babel}
\usepackage{actuarialangle, actuarialsymbol}


\begin{document}

Define the Cramer-Lundberg model

If one would like to consider instead \(t \in \mathbb{N}_0\) (discrete time)
continuous time (i.e. \(t \in \mathbb{R}_+ = [0, \infty)\)),
one can use the Cramér-Lundberg model.

First step: 
Let \((N_t)_{t \in \mathbb{R}_+}\) be Poisson process with intensity \(\lambda v > 0\)
\(N_t \coloneqq \) total number of claims incurred up to time \(t\)

Definition (Poisson Process):
Counting process \((N_t)_{t \in \mathbb{R}_+}\) is Poisson process with intensity \(\lambda v > 0\) if
\(N_0 = 0\)
independent and stationary increments
\( \mathbb{P}_{N_t}(n) = e^{-\lambda vt} \frac{(\lambda v t)^n}{n!}, n \in \mathbb{N}_0\) 
(Poisson-distribution)

Second step:
Total premia up to \( t \in \mathbb{R}_{+}\) is \( \beta vt \) for (\( \beta > 0 \) fixed constant)
Assume that \( (S_n)_{n \in \mathbb{N}} \) is a seq. of i.i.d random variables
and \( \sigma(S_n \mid n \in \mathbb{N}) \) is independent of \( \sigma(N_t \mid t \in \mathbb{R}_{+}) \),
then we define the total claim size up to time \( t \) to be \( \sum_{n=1}^{N_t} S_n \).

Summarizing
The Cramer-Lundberg surplus process for initial capital \( c_0 \geq 0 \) is
\( \tilde{C}_t = c_0 + \beta vt - \sum_{n=1}^{N_t} S_n \) for \( t \in \mathbb{R}_{+} \)
Its ruin time is
\( \tilde{\tau} = \inf\{t \in \mathbb{R}_{+}\mid \tilde{C}_t < 0\}\)



Make some observations in regard to the Cramer-Lundberg model 

\(\tilde{\tau}\) can only occur at "jump times" of \(N\)

Define
\(W_u := \inf \{t \in \mathbb{R}_+ | N_t = u\} - \inf \{t \in \mathbb{R}_+ | N_t = u - 1\}, u \in \mathbb{N}\)
\(W_u \hat{=}\) time \(N\) spends at value \(u - 1\)

\(N\) Poisson process of intensity \(\lambda v \implies (W_u)_{u \in \mathbb{N}}\) i.i.d. Exp(\(\lambda v\))-random variables
\(N\) ind. of \((S_n)_{n \in \mathbb{N}} \implies (W_u)_{u \in \mathbb{N}}\) ind. of \((S_n)_{n \in \mathbb{N}}\)
\(V_n := \sum_{u = 1}^n W_u, n \in \mathbb{N}\)
\(\rightarrow\) in order to study \(\tilde{\tau}\), suffices to study the sequence
\(C_n := \tilde{C}_{V_n} = c_0 + \beta v V_n - \sum_{u = 1}^n S_u = c_0 + \sum_{u = 1}^n (\beta v W_u - S_u), n \in \mathbb{N}\)


Explain the idea behind premium principles

Paradigm: 
The premium should depend on random claim size \( S \)
and the NPC.

Let \( S \in L^1 \) be a random claim
\( \mathbb{E}[S] \) is called the pure risk premium
The difference
\( \pi(S) - \mathbb{E}[S] \) is called the safety margin.
Hence
\( \text{Premium} = \text{pure risk premium} + \text{safety margin} \).

There are different principles to determine the safety margin, for example

Let \( S \geq 0 \) be a random claim with finite first moment and \( \alpha > 0 \). 
Then \( \pi(S) \coloneqq (1 + \alpha)\mathbb{E}[S] \)
is called expected value principle.

The (NPC) is satisfied \(S ≥ 0\) in this case,
nevertheless it is problematic because insensitive to randomness of \( S \).
For example \( k > 0, S \geq 0 \) any random variable with \( \mathbb{E}[S] = k \implies \pi(S) = (1 + \alpha)k \)
In particular: \( S < \pi(S) \) possible! (e.g. S ≡ k)

Define the variance loading principle for premia

Suppose the claim \( S \geq 0 \) satisfies \( \mathbb{E}[S^2] < \infty\). Let \( \alpha > 0 \)
then \( \pi(S) \coloneqq \mathbb{E}[S] + \alpha \text{Var}[S] \)
is called variance loading principle.
As an example consider,
\(\Omega = (0, 1), \mathcal{F} = \mathbb{B}((0, 1)), \mathbb{P} =\) Lebesgue measure
\(A := (0, \frac{1}{2}), k > 0, 0 < \beta < \gamma < k\)
\(S_1 := k, S_2 := (k - \beta)\mathbf{1}_A + (k + \beta)\mathbf{1}_{A^c}, S_3 := (k - \gamma)\mathbf{1}_A + (k + \gamma)\mathbf{1}_{A^c}\)
\(\implies \pi(S_1) = k < \pi(S_2) = k + \alpha \beta^2 < \pi(S_3) = k + \alpha \gamma^2\)

This is however tricky to use directly in reality, namely that \( \text{Var}[S] \) usually involves some kind
of square term and the square of € has no direct meaning.


Suppose claim \(S \ge 0\) satisfies \(\mathbb{E}[S^2] < \infty\). Let \(\alpha > 0\).
\(\pi(S) := \mathbb{E}[S] + \alpha \sqrt{\text{var}(S)} = \mathbb{E}[S] + \alpha \cdot sd(S)\)
is called standard deviation loading principle (sd)-loading principle.

Remark:
If \(\mathbb{E}[S] > 0, \pi =\) sd-loading principle, then
\(\pi(S) = \mathbb{E}[S] \left( 1 + \frac{\sqrt{\text{var}(S)}}{\mathbb{E}[S]} \right).\)

The previous examples then becomes:
\( π(S1) = k < π(S2) = k + αβ < π(S3) = k + αγ \).
Quantity \(Vco(S) := \frac{\sqrt{\text{var}(S)}}{\mathbb{E}[S]}\) called coefficient of variation. Sd-loading clearly risk based.
Sd-loading principle is positively homogeneous in S.


Roughly summarize Utility Theory


Key idea (in a nutshell): Premia are derived from preferences of financial agents
Leads to utility theory
UT measures "happiness index" of economic agents with (random) endowment
\(\Omega, \mathcal{F}, \mathbb{P})\): probability space
\(L^1 = L^1(\Omega, \mathcal{F}, \mathbb{P}):\) space of (equivalence classes of) integrable random variables
\(\mathcal{X} \subset L^1, I \subset \mathbb{R}\) interval s.t.
\(\forall X \in \mathcal{X}: \mathbb{P}(X \in I) = 1\)
Preferences of agent captured by binary relation \(\succeq \text{ on } \mathcal{X}\)
\(X \succeq Y \equiv\), "\(Y\) is weakly preferred to \( X \)"
\(X \sim Y : \iff X \succeq Y\) and \(Y \succeq X \equiv\) "indifference between \(X\) and \(Y\)"

Numerical representation: Function 
\(\mathcal{U}: \mathcal{X} \rightarrow [-\infty, \infty] = \mathbb{R} \cup \{\pm \infty\}\) s.t. 
\(X \preceq Y \iff \mathcal{U}(X) \le \mathcal{U}(Y)\)

Expected utility: We will focus on "von Neumann-Morgenstern preferences"
i.e., \(\exists u: I \rightarrow \mathbb{R}\) such that \(\mathcal{U}(X) := \mathbb{E}[u(X)], X \in \mathcal{X}\) 
represents \(\preceq\).

Some further assumptions are:
\( I^{°} \) (interior of \(I\)) is nonempty
Utility function \(u: I \rightarrow \mathbb{R}\) has the following properties:

1. strict monotonicity: \(x, y \in I, x < y \implies u(x) < u(y)\) ("more is better than less")
2. risk aversion: \(u\) is strictly concave, i.e., 
    \(x, y \in I, x \ne y, \lambda \in (0, 1) \implies u(\lambda x + (1 - \lambda) y) > \lambda u(x) + (1 - \lambda) u(y)\)

Interpret \(X \in \mathcal{X}\) as gain

\(X, Y \in \mathcal{X}, X \preceq Y\) & \(X \ne Y \implies \mathbb{E}[u(X)] < \mathbb{E}[u(Y)]\) (\(u\) strictly increasing!)
\(\implies Y\) strictly preferred to \(X\)
\(u\) strictly concave
Jensen inequality \(\implies \mathbb{E}[u(X)] \le u(\mathbb{E}[X]) = \mathbb{E}[u(\mathbb{E}[X])]\)
Preceding inequality strict if \(X\) not deterministic
\(\implies\) agent risk averse in that certainty (deterministic \(\mathbb{E}[X]\)) always preferred to randomness (\(X\)).


Define utility indifference prices and state their uniqueness

\( c_0 \in I \) initial capital
\( S \in L^1 \)
Any number \( \pi \in \mathbb{R} \), s.t. \(\mathbb{P}(c_0 + \pi - S \in I) = 1\) is called a \(\textbf{price}\) of \( S \).
The \(\textbf{utility indifference price}\) of  \(S\) is a price \(\pi(S)\) of \(S\) s.t.
\(u(c_0) = \mathbb{E}[u(c_0 + \pi(S) - S)]\)
(provided such a solution exists).

From a utility theory perspective this is equivalent to \( c_0 \preceq c_0 - Y + \pi \).

Proposition:
Suppose \( S \in L^1 \) is such that utility indifference price \( π(S) \) exists.
1. \( π(S) \) is unique.
2. If \( S \) is not deterministic and \( \mathbb{E}[S] \) is a price of \( S \), then \( π(S) > E[S] \) (the (NPC)).
Proof:
1. \( u \) is strictly increasing, thus \( u(c_0) = \mathbb{E}[u(c_0 + \pi_1 - S)] = \mathbb{E}[u(c_0 + \pi_0 - S)] \)
is not possible for \( \pi_1 \neq \pi_0 \).


Define CARA utilities and state their universal property

\(I = \mathbb{R}\)
\(\alpha > 0\) fixed constant
\(u: \mathbb{R} \rightarrow \mathbb{R}\) defined by \(u(x) = 1 - \frac{1}{\alpha} e^{-\alpha x}\)
\(u\) called CARA utility function (CARA = "constant absolute risk aversion")

As an example:
Suppose \(S \sim N(\mu, \sigma^2)\)
Need to solve:

\(1 - \frac{1}{\alpha} e^{-\alpha c_0} = \mathbb{E} \left[ 1 - \frac{1}{\alpha} e^{-\alpha (c_0 + \pi(S) - S)} \right]\)
\(1 = e^{-\alpha \pi(S)} \mathbb{E} \left[ e^{\alpha S} \right]\)
\(e^{\alpha \pi(S)} = \exp \left( \alpha \mu + \frac{1}{2} \alpha^2 \sigma^2 \right)\)
\(\pi(S) = \mu + \frac{1}{2} \alpha \sigma^2 = \mathbb{E}[S] + \frac{1}{2} \alpha \text{var}(S).\)

Observe: \(\pi(S)\) not dependent on \(c_0\)!

The universal property of CARA utilities is 
Suppose:
* \(I = \mathbb{R}\)
* \(u \in C^2(\mathbb{R})\)
* \(\forall p \in (0, 1), \exists A_p \in \mathcal{F}: \mathbb{P}(A_p) = p.\)
(atomless spaces satisfy this for example).

Then the following are equivalent:
For all \(S \in L^1\) such that \(\pi(S)\) exists, the latter does not depend on \(c_0 \in \mathbb{R}\).
\(\exists a \in \mathbb{R}, b, c > 0\): \(u(x) = a - be^{-cx}, x \in \mathbb{R}\).
As one can see CARA utilities can be transformed into this shape and vice versa.


Define absolute/relative risk aversion of utility functions
and relate it to CARA

\(u \in C^2(\mathbb{R})\) utility function. Absolute/relative risk aversion of \(u\):
\(\rho_{ARA}, \rho_{RRA}: \mathbb{R} \rightarrow \mathbb{R}\)
\(\rho_{ARA}(x) := -\frac{u''(x)}{u'(x)}, \rho_{RRA}(x) := -\frac{xu''(x)}{u'(x)}.\)

Example:
\(\alpha > 0, u_\alpha(x) := 1 - \frac{1}{\alpha} e^{-\alpha x}, x \in \mathbb{R}\)
\(u_\alpha'(x) := e^{-\alpha x}, u_\alpha''(x) := -\alpha e^{-\alpha x}, x \in \mathbb{R}\)

Hence
\(\rho_{ARA}(x) = -\frac{-\alpha e^{-\alpha x}}{e^{-\alpha x}} = \alpha\)
Explains "CARA = constant absolute risk aversion"

This motivatives also the following family
Let \( u_\gamma(x) = (1 - \gamma)^{-1}x^{1 - \gamma} \) for \( \gamma \neq 1 \) and otherwise \( u_\gamma(x) = \log(x) \),
then \( \rho^{RRA}(x) = - \frac{x(-\gamma x^{-\gamma - 1}}{x^{-\gamma}} = \gamma\),
which are called CRRA utilities (constant relative risk aversion).


Define risk aversion and its relation to premia

\(I \subset \mathbb{R}\) open interval
\(u, v: I \rightarrow \mathbb{R}\) utility functions
\(u\) more risk averse than \(v\) if
\(\forall X \in L^1: \mathbb{P}(X \in I) = 1, \mathbb{E}[u(X)], \mathbb{E}[v(X)] \in \mathbb{R}\)
implies
\(u^{-1}(\mathbb{E}[u(X)]) \le v^{-1}(\mathbb{E}[v(X)]).\)

\(u^{-1}(\mathbb{E}[u(X)])\) and \(v^{-1}(\mathbb{E}[v(X)])\) are the respective certainty equivalents of \(X\):
\(c := u^{-1}(\mathbb{E}[u(X)])\) satisfies \(u(c) = \mathbb{E}[u(X)]\)
and similarly for \(v\).

If \(u\) more risk averse than \(v\) and
\(c_0 \in I, S \in L^1\) s.t. \(\pi^u(S)\) and \(\pi^v(S)\) exist then \(\pi^v(S) \le \pi^u(S)\).

Motivate and define the Esscher premium principle

Assume \(S \in L^1\) and \(\mathbb{E}[e^{\alpha S}] < \infty\) for some \(\alpha > 0\)
and let \(F \hat{=}\) be the CDF of \(S\).
Classical actuarial reasoning: "bad events" (loss events) should count more than good events,
so they need more weight. This can be done by modifying \( F \), which is precisely what the
Escher premium does.


Fix \(\alpha > 0\) s.t. \(M_S(\alpha) := \int_{\mathbb{R}} e^{\alpha x} dF(x) = \mathbb{E}[e^{\alpha S}] < \infty\). 
The CDF
\(F_\alpha: \begin{cases} \mathbb{R} \rightarrow [0, 1], \\ x \mapsto M_S(\alpha)^{-1} \int_{-\infty}^x e^{\alpha s} dF(s) \end{cases}\)
is the Esscher transform of \(F\).
Assume there is \(r_0 > 0\) s.t. \(M_S(r) < \infty, r \in (-r_0, r_0)\). For \(\alpha \in (0, r_0)\)
\(\pi_\alpha(S) := \int_{\mathbb{R}} s dF_\alpha(s) = \frac{\mathbb{E}[Se^{\alpha S}]}{\mathbb{E}[e^{\alpha S}]}\)
is called Esscher premium of \(S\).

We see that \( \pi_\alpha(S) \) emphasises losses exponentially
and it can be shown that \( \pi_\alpha(S) \) is finite.

In particular, the Esscher premium satisfies the NPC: 
\( \mathbb{E}[S] \leq \pi_\alpha(S) \)
with more effort the strict inequality holds for non deterministic \( S \in L^1 \).
Though the Esscher premium can only be computed for light-tailed claims!

In the proofs of these results, one requires the function \( h : (-r_0, r_{0}) \to \mathbb{R}, r \mapsto \log(M_S(r)) = \mathbb{E}[e^{\alpha S}]\),
which is convex and the modified version \( \frac{1}{\alpha}h(\alpha) \) is referred to as entropic risk measure.
\end{document}
